\chapter*{ВВЕДЕНИЕ}
\addcontentsline{toc}{chapter}{ВВЕДЕНИЕ}

\section*{Актуальность темы}

Распределённые вычислительные системы составляют основу цифровой инфраструктуры \cite{google_sre_book}. Рост сложности и масштаба таких систем приводит к увеличению частоты инцидентов и сложности их диагностики.

Статистика показывает, что среднее время восстановления (MTTR) после инцидента в крупных распределённых системах составляет от 30 минут до нескольких часов \cite{netflix_chaos}. При этом до 70\% времени восстановления тратится на диагностику проблемы, а не на её устранение \cite{pagerduty_incident_response}. Это генерирует прямые экономические потери: час простоя критического сервиса обходится в \$100\,000--\$500\,000 \cite{amazon_outage_cost}.

Традиционные подходы к мониторингу, основанные на метриках и пороговых значениях, не справляются с растущей сложностью современных систем. Они дают много ложных срабатываний, не учитывают взаимосвязи между компонентами системы и не способны предсказывать деградацию до возникновения критических проблем \cite{charity_majors_observability}.

Концепция наблюдаемости (observability) предлагает новый подход к пониманию состояния системы через анализ метрик, логов и трейсов. Однако современные инструменты наблюдаемости требуют от инженеров ручной интерпретации данных \cite{charity_majors_observability}.

\begin{figure}[H]
\centering
\begin{tikzpicture}[scale=0.8]
\begin{axis}[
    xlabel={Год},
    ylabel={Количество инцидентов (тыс.)},
    title={Рост количества инцидентов в распределённых системах},
    grid=major,
    legend pos=north west,
]
\addplot[blue,thick] coordinates {
    (2015, 50) (2016, 65) (2017, 85) (2018, 110) (2019, 140) (2020, 180) (2021, 220) (2022, 270) (2023, 320)
};
\addplot[red,thick,dashed] coordinates {
    (2015, 30) (2016, 35) (2017, 40) (2018, 45) (2019, 50) (2020, 55) (2021, 60) (2022, 65) (2023, 70)
};
\legend{Общее количество инцидентов, Количество критических инцидентов}
\end{axis}
\end{tikzpicture}
\caption{Рост количества инцидентов в распределённых системах за последние годы \cite{aiops_survey,pagerduty_incident_response}}
\label{fig:incidents_growth}
\end{figure}

\begin{table}[H]
\centering
\caption{Статистика времени восстановления после инцидентов}
\label{tab:recovery_time_stats}
\begin{tabular}{|p{5cm}|c|c|c|}
\hline
\textbf{Тип инцидента} & \textbf{MTTR, мин} & \textbf{Диагностика, \%} & \textbf{Восстановление, \%} \\
\hline
Отказ сервиса & 45 & 65 & 35 \\
\hline
Деградация производительности & 120 & 70 & 30 \\
\hline
Утечка памяти & 180 & 55 & 45 \\
\hline
Каскадный отказ & 240 & 60 & 40 \\
\hline
\end{tabular}
\end{table}

Технологии машинного обучения и больших языковых моделей позволяют автоматизировать диагностику и восстановление систем. Интеллектуальные автономные агенты анализируют потоки метрик, логов и трейсов, выявляют паттерны аномалий, строят причинно-следственные связи и принимают решения о восстановлении \cite{aiops_survey}.

\section*{Состояние разработанности проблемы}

Традиционные системы мониторинга, такие как Nagios, Zabbix, Prometheus, собирают метрики и генерируют алерты, но их пороговые значения задаются вручную, а сложные паттерны аномалий не выявляются автоматически \cite{prometheus_docs}. Современные платформы наблюдаемости, включающие Grafana, ELK Stack, Jaeger, предоставляют более продвинутые возможности анализа, но всё ещё требуют значительного участия инженеров в интерпретации данных \cite{grafana_docs,charity_majors_observability}.

Подходы на основе машинного обучения для обнаружения аномалий, такие как методы временных рядов, изолирующие леса, нейронные сети, показывают хорошие результаты в задачах детектирования, но часто не учитывают причинно-следственные связи между компонентами системы \cite{hastie_elements,bishop_pattern_recognition}.

Архитектуры автономных агентов, включая реактивные системы и BDI-архитектуры, демонстрируют потенциал для автоматизации операций, но существующие реализации часто ограничены простыми правилами и не используют современные достижения в области больших языковых моделей для анализа логов и генерации гипотез \cite{aiops_survey}.

Однако комплексное решение, объединяющее обнаружение аномалий, причинно-следственный анализ и безопасное самовосстановление в единую архитектуру интеллектуальных агентов, в литературе не описано \cite{aiops_survey}.

\section*{Проблема изучения темы}

Существует противоречие между:
\begin{itemize}
    \item Ростом сложности распределённых систем и необходимостью быстрого обнаружения и устранения проблем;
    \item Ограниченными возможностями традиционных систем мониторинга и потребностью в интеллектуальной диагностике;
    \item Большими объёмами данных наблюдаемости и необходимостью их эффективного анализа для принятия решений;
    \item Требованиями к безопасности операций восстановления и необходимостью автоматизации процессов.
\end{itemize}

\section*{Цель}

Разработать архитектуру и методы построения интеллектуальных автономных агентов, которые диагностируют, мониторят и восстанавливают распределённые вычислительные системы, опираясь на комплексный анализ метрик, логов и трейсов.

\section*{Задачи исследования}

Необходимо решить следующие задачи:

\begin{enumerate}
    \item Провести анализ теоретических основ распределённых вычислительных систем, их архитектур и принципов построения устойчивых систем.

    \item Исследовать концепцию наблюдаемости как фундамента диагностики распределённых систем, проанализировать современные подходы к сбору и анализу метрик, логов и трейсов.

    \item Провести сравнительный анализ backend-подходов к построению устойчивых систем: синхронных и асинхронных архитектур, REST и gRPC, event-driven систем, CQRS, принципов идемпотентности и паттернов устойчивости.

    \item Изучить инженерные парадигмы DevOps и SRE, методы автоматизации развёртывания, стратегии самовосстановления в Kubernetes, ограничения современных систем мониторинга.

    \item Исследовать методы обнаружения аномалий и прогнозирования деградации: статистические методы, методы машинного обучения для временных рядов, применение LLM для анализа логов.

    \item Проанализировать архитектуры интеллектуальных автономных агентов: реактивные архитектуры, BDI-архитектуры, rule-based, ML-based и LLM-based подходы, жизненный цикл агента.

    \item Разработать архитектуру интеллектуального автономного агента для диагностики и самовосстановления, определить требования, архитектурные принципы, модули системы, механизмы оркестрации и безопасного исполнения решений.

    \item Провести сравнительный анализ инженерных парадигм: монолит и микросервисы, push- и pull-модели, event-driven и request-driven архитектуры, реактивные и акторные модели, rule-based и ML-based диагностика.

    \item Реализовать прототип системы интеллектуальных автономных агентов, разработать структуру модулей, схемы взаимодействия, параметры конфигурации, моделирование сценариев самовосстановления.

    \item Провести экспериментальное исследование эффективности разработанного подхода: разработать методику экспериментов, определить сценарии тестирования, метрики оценки (MTTR, latency, precision, recall), провести анализ результатов.

    \item Оценить практическую значимость разработанного решения: определить области применения в индустрии, описать возможности интеграции в реальные системы.
\end{enumerate}

\section*{Объект исследования}

Объект исследования --- распределённые вычислительные системы и процессы их диагностики, мониторинга и восстановления.

\section*{Предмет исследования}

Предмет исследования --- методы и алгоритмы построения интеллектуальных автономных агентов для автоматической диагностики и самовосстановления распределённых систем на основе анализа данных наблюдаемости.

\section*{Новизна работы}

Научная новизна:

\begin{enumerate}
    \item Разработана комплексная архитектура интеллектуальных автономных агентов, объединяющая методы машинного обучения для обнаружения аномалий, причинно-следственный анализ инцидентов и безопасные механизмы самовосстановления.

    \item Предложен новый подход к построению причинно-следственных графов инцидентов на основе комплексного анализа метрик, логов и трейсов, позволяющий повысить точность диагностики корневых причин проблем.

    \item Разработана методология интеграции больших языковых моделей для анализа логов и автоматической генерации гипотез о причинах инцидентов.

    \item Предложен механизм безопасного самовосстановления с использованием fail-safe принципов и валидации действий агентов перед их выполнением.

    \item Разработана система метрик для оценки эффективности интеллектуальных автономных агентов, учитывающая не только точность обнаружения аномалий, но и безопасность и эффективность действий по восстановлению.
\end{enumerate}

\section*{Научная гипотеза}

Использование интеллектуальных автономных агентов, объединяющих методы машинного обучения для обнаружения аномалий, причинно-следственный анализ инцидентов и безопасные механизмы самовосстановления, снижает MTTR и повышает доступность распределённых вычислительных систем по сравнению с пороговым мониторингом и ручной диагностикой.

Формально, пусть \(MTTR\) обозначает среднее время восстановления после инцидента, а \(R\) — интегральную метрику надёжности системы (например, вероятность безотказной работы на заданном интервале времени). Гипотеза утверждает: существует конфигурация агентов и процедура их интеграции, при которой выполняются неравенства
\begin{equation}
  MTTR_{\text{agent}} < MTTR_{\text{baseline}}, \quad
  R_{\text{agent}} > R_{\text{baseline}},
\label{eq:hypothesis_ineq}
\end{equation}
где индекс \(\text{baseline}\) соответствует существующей системе без интеллектуальных агентов, а \(\text{agent}\) — системе с внедрёнными автономными агентами диагностики и самовосстановления.

\section*{Методологическая база исследования}

Методы исследования:

\begin{enumerate}
    \item \textbf{Теоретический анализ} — изучение научной литературы, стандартов и документации по распределённым системам, наблюдаемости, машинному обучению и архитектурам агентов.

    \item \textbf{Сравнительный анализ} — сопоставление различных подходов к построению устойчивых систем, методов обнаружения аномалий, архитектур агентов.

    \item \textbf{Математическое моделирование} — разработка моделей для анализа временных рядов метрик, построения причинно-следственных графов, принятия решений агентами.

    \item \textbf{Проектирование архитектуры} — разработка архитектуры интеллектуальных автономных агентов с использованием принципов модульности, расширяемости и безопасности.

    \item \textbf{Экспериментальное моделирование} — создание прототипа системы и проведение экспериментов на синтетических и реальных данных.

    \item \textbf{Статистический анализ} — оценка эффективности разработанного подхода через метрики точности, полноты, времени восстановления.
\end{enumerate}

\section*{Положения, выносимые на защиту}

На защиту выносятся следующие положения:

\begin{enumerate}
    \item Архитектура интеллектуальных автономных агентов для диагностики и самовосстановления распределённых систем, интегрирующая LSTM-детектор аномалий, причинно-следственные графы инцидентов и механизм безопасного исполнения с формальной оценкой риска.

    \item Метод построения причинно-следственных графов инцидентов на основе комплексного анализа метрик, логов и трейсов, обеспечивающий повышение точности диагностики корневых причин проблем по сравнению с существующими подходами.

    \item Методология интеграции больших языковых моделей для анализа логов и автоматической генерации гипотез о причинах инцидентов, позволяющая сократить время диагностики.

    \item Механизм безопасного самовосстановления с использованием fail-safe принципов и валидации действий агентов перед их выполнением, обеспечивающий надёжность операций восстановления.

    \item Система метрик для оценки эффективности интеллектуальных автономных агентов, учитывающая не только точность обнаружения аномалий, но и безопасность и эффективность действий по восстановлению.
\end{enumerate}

\section*{Практическое значение}

Практическая значимость работы:

\begin{enumerate}
    \item Архитектура и алгоритмы применимы при проектировании систем мониторинга и автоматизации операций в IT-компаниях и облачных провайдерах.

    \item Автоматизация диагностики и восстановления сокращает операционные затраты и время простоя.

    \item Предложенные подходы могут быть интегрированы в существующие платформы наблюдаемости (Prometheus, Grafana, OpenTelemetry) для расширения их функциональности \cite{prometheus_docs,grafana_docs,opentelemetry_docs}.

    \item Разработанные методы могут быть использованы при внедрении практик SRE (Site Reliability Engineering) в организациях \cite{google_sre_book,sre_principles}.

    \item Результаты работы могут быть применены при проектировании систем для критически важных приложений: финансовых систем, систем здравоохранения, промышленной автоматизации \cite{iso_25010,amazon_outage_cost}.
\end{enumerate}

Ожидаемые улучшения ключевых метрик суммированы в таблице~\ref{tab:expected_improvements}.

\begin{table}[H]
\centering
\caption{Ожидаемые улучшения метрик при использовании интеллектуальных агентов}
\label{tab:expected_improvements}
\begin{tabular}{|p{5cm}|c|c|c|}
\hline
\textbf{Метрика} & \textbf{Текущее} & \textbf{Ожидаемое} & \textbf{Улучшение} \\
\hline
MTTR, мин & 79 & 44 & -44\% \\
\hline
Точность обнаружения, \% & 72 & 92 & +28\% \\
\hline
Ложные срабатывания, \% & 28 & 8 & -71\% \\
\hline
\end{tabular}
\end{table}

Для каждой метрики
\[
  m \in \{
    \text{MTTR},
    \text{precision},
    \text{false\_positive\_rate},
    \text{diagnosis\_time}
  \}
\]
относительное улучшение определяется как
\begin{equation}
  I_m = \frac{m_{\text{baseline}} - m_{\text{agent}}}{m_{\text{baseline}}} \cdot 100\%,
\end{equation}
где \(m_{\text{baseline}}\) — значение метрики в существующей системе, \(m_{\text{agent}}\) — значение метрики при использовании интеллектуальных агентов. Для метрик, интерпретируемых по принципу «чем больше, тем лучше» (например, точность обнаружения аномалий), используется модифицированное определение:
\begin{equation}
  I_m = \frac{m_{\text{agent}} - m_{\text{baseline}}}{m_{\text{baseline}}} \cdot 100\%.
\end{equation}


\section*{Структура работы}

ВКР состоит из введения, четырёх разделов, заключения, списка использованных источников, приложений и актов внедрения.

\begin{description}
\item[\textbf{Введение}] обосновывает актуальность, формулирует цель, задачи, новизну, гипотезу и положения, выносимые на защиту.

\item[\textbf{Раздел 1}] <<Обзор литературных источников по теме диплома, аналогов, прототипов>> содержит анализ теоретических основ распределённых вычислительных систем, концепции наблюдаемости, backend-подходов к устойчивым системам, инженерных парадигм DevOps и SRE, методов обнаружения аномалий и архитектур интеллектуальных автономных агентов.

\item[\textbf{Раздел 2}] <<Методы создания системы, реализация>> представляет архитектуру проектируемой системы, сравнительный анализ инженерных парадигм и описание реализации прототипа.

\item[\textbf{Раздел 3}] <<Результаты работы>> содержит описание экспериментального исследования: методику экспериментов, сценарии тестирования, метрики оценки и результаты экспериментов.

\item[\textbf{Раздел 4}] <<Экономическая эффективность разработки>> оценивает практическую значимость разработанного решения, области применения в индустрии, возможности интеграции в реальные системы и экономическую оценку эффекта.

\item[\textbf{Заключение}] содержит основные выводы, результаты работы и перспективы развития.
\end{description}